\section{Introducción}
\label{ch:introduccion}
%Contextualización del tema%
En el ámbito de las operaciones tecnológicas corporativas, la gestión eficiente de tareas y la organización visual de la información son fundamentales para garantizar la continuidad de los servicios \cite{manza-diaz_gestion_2025}. En la empresa Honne S.A. de C.V., específicamente dentro de su Centro de Operación y Excelencia Multicloud (CCoE), el equipo lleva a cabo actividades críticas de alta exigencia para clientes a nivel internacional, abarcando desde el soporte técnico y el monitoreo preventivo de infraestructuras (NOC/SOC), hasta la validación de integraciones de software \cite{honne_oficial_2026}. Tradicionalmente, la delegación, asignación y seguimiento de estas responsabilidades se han administrado mediante el uso de un documento en Excel, el cual funciona como un manual operativo y agendario central. Sin embargo, la gestión de entornos tecnológicos complejos a través de hojas de cálculo plantea retos significativos en la interacción diaria de los empleados, lo que hace evidente la necesidad de evolucionar hacia plataformas que prioricen una experiencia de usuario (UX) fluida y un diseño de interfaz (UI) estructurado para soportar altas cargas operativas \cite{gupta_impact_2026}.


%Problema de investigación%
El problema central de esta investigación radica en la ineficiencia y la saturación visual que genera el uso de Microsoft Excel como herramienta principal para la delegación de tareas en el Centro de Operación y Excelencia Multicloud (CCoE) de Honne S.A. de C.V. Actualmente, el personal se enfrenta a una interfaz estática y desorganizada, donde la información crítica —como el filtrado de actividades (Programadas, Internas, Tibco), el calendario de turnos y las notificaciones de ausencias— se encuentra aglomerada y sin diferenciación por perfiles de usuario (Administrador vs. Empleado). Esta falta de organización visual provoca una navegación confusa, vuelve más lentos los procesos diarios y eleva el riesgo de errores humanos en un entorno que exige respuestas rápidas y alta disponibilidad. Por lo tanto, el presente estudio se delimita espacial y poblacionalmente al personal operativo y administrativo del CCoE, buscando dar respuesta a la siguiente interrogante: ¿De qué manera el diseño de una interfaz web estructurada bajo principios de usabilidad (UX/UI) puede resolver la saturación de información y optimizar el seguimiento de las métricas operativas en el entorno corporativo?

%Justificación, Objetivos, Metodología y Estructura (luego lo extiendo)%
El objetivo principal de este proyecto es diseñar una interfaz gráfica intuitiva y centrada en el usuario que reemplace el actual manual en Excel, reduciendo el tiempo de interacción y mejorando la eficiencia del personal del CCoE. Esta propuesta se justifica por la necesidad de mitigar errores humanos y agilizar la visualización de tareas operativas mediante metodologías de diseño UX/UI y arquitectura de la información. A lo largo de este documento se detallan todos los aspectos relevantes de la investigación, abordando desde los objetivos establecidos y la justificación del diseño, hasta la metodología aplicada, los alcances, limitaciones y los resultados obtenidos en la propuesta visual.

\subsection{Antecedentes}
\label{sec:antecedentes}
En los últimos años, la gestión operativa dentro de las organizaciones tecnológicas ha experimentado una transición acelerada, pasando de procesos administrativos manuales a la adopción de plataformas digitales centralizadas. Esta evolución ha permitido a las empresas escalar el alcance de sus servicios, pero al mismo tiempo ha incrementado la complejidad de la información que los empleados deben procesar diariamente para llevar a cabo sus funciones [4].

Históricamente, la organización interna, el seguimiento de tareas y la comunicación de turnos se apoyaba fuertemente en herramientas ofimáticas genéricas, como hojas de cálculo o documentos de texto compartidos. Sin embargo, en entornos corporativos de alta exigencia, depender de interfaces estáticas y no especializadas resulta ineficiente. La falta de jerarquía visual y la aglomeración de datos se traducen en una alta carga cognitiva para el usuario, lo que frecuentemente deriva en lentitud operativa, pérdida de información crítica y errores humanos [5].

Para mitigar estos riesgos, la industria del desarrollo de software ha comenzado a priorizar el Diseño de Experiencia de Usuario (UX) y de Interfaces de Usuario (UI) no solo en productos comerciales, sino en las herramientas corporativas internas. La implementación de gestores web estructurados, con navegación intuitiva y vistas personalizadas según el rol del empleado, se ha vuelto un estándar para reducir la fricción en el trabajo diario. El presente proyecto surge ante la necesidad de modernizar las herramientas de gestión del personal, garantizando un entorno digital que facilite la toma de decisiones y optimice el rendimiento operativo [5].

\subsubsection{Antecedentes de la empresa}
\label{sec:antecedentes_empresa}

Honne S.A. de C.V. es una empresa especializada en transformación tecnológica que acompaña a las organizaciones en la evolución de su operación, desde la estrategia hasta la ejecución. A través de sus capacidades avanzadas en tecnologías de la nube (\textit{cloud}), gestión de datos e Inteligencia Artificial, la empresa ayuda a construir entornos más eficientes, escalables y preparados para el crecimiento de diversos sectores industriales\cite{honne_oficial_2026}.

Su misión institucional es ``Crear valor para nuestros clientes a través de la Innovación y las Tecnologías de la Información, utilizando nuestro portafolio de servicios, experiencia y capacidades; para resolver sus desafíos actuales y futuros''. Asimismo, su visión se enfoca en ``Ser la empresa preferida a nivel mundial para la gestión de tecnologías de nube pública y la creación de servicios de próxima generación, de forma segura y alineada con las mejores prácticas''\cite{honne_oficial_2026}..

Para garantizar la calidad y seguridad en cada proyecto, Honne opera bajo rigurosos estándares internacionales y certificaciones clave en la industria, tales como ISO 9001, ISO 20000, ISO 27001, ITIL, CMMI y TOGAF. Su cultura laboral se rige por la filosofía \textit{One Team\textregistered}, la cual promueve la integración total con cada cliente asumiendo un alto nivel de compromiso, empatía y sentido de urgencia frente a los desafíos tecnológicos\cite{honne_oficial_2026}..

El desarrollo del presente proyecto se enmarca dentro de su Centro de Operación y Excelencia Multicloud (CCoE), con sede en Ciudad Victoria, Tamaulipas. Desde este núcleo tecnológico, la empresa gestiona y opera entornos de TI para clientes en México, Estados Unidos, Europa y Latinoamérica. Es en este entorno operativo donde se ejecutan las labores de soporte técnico, el monitoreo preventivo de infraestructuras (NOC/SOC) y la validación de integraciones, asegurando así la continuidad, eficiencia y calidad ininterrumpida de los servicios empresariales\cite{honne_oficial_2026}..

\subsubsection{Antecedentes teóricos y Trabajos Relacionados}
\label{sec:antecedentes_teoricos}

\textbf{Centros de Operaciones (NOC/SOC) y Monitoreo Proactivo} \\ 
Un Centro de Operaciones de Red (NOC) es una ubicación centralizada donde administradores y técnicos de TI supervisan, monitorean y mantienen las redes y telecomunicaciones de una empresa de forma ininterrumpida. La evolución de los NOC modernos requiere el uso de herramientas de monitoreo de infraestructura para vigilar métricas críticas y prevenir fallas. 

En este contexto, Arias (2025) desarrolló el proyecto SIMAI, proponiendo una arquitectura basada en Microsoft Azure para abordar las limitaciones del monitoreo tradicional, buscando centralizar alertas y reducir la intervención manual \cite{Arias_Pro_2025}. Sin embargo, dicho trabajo se limitó a un diseño conceptual sin implementación práctica \cite{Arias_Pro_2025}. A diferencia de esa propuesta, el \textbf{\NombreProyecto} supera esa limitante al ejecutar validaciones en un entorno de producción real, aplicando las metodologías directamente sobre cuentas y aplicativos de clientes corporativos activos.

\textbf{Gestión de Servicios de TI (ITSM) y SLA} \\ 
La Gestión de Servicios de TI (ITSM) se refiere al enfoque estratégico para diseñar, entregar, gestionar y mejorar la forma en que las empresas utilizan la tecnología de la información. La ejecución eficiente de estas operaciones requiere alinearse a marcos de referencia globales como ITIL v4, el cual proporciona lineamientos estructurados para la gestión del ciclo de vida de los incidentes, la mejora continua y la entrega de valor al negocio \cite{Arias_Pro_2025}. 

Un componente crítico dentro de la arquitectura ITSM es el estricto cumplimiento de los Acuerdos de Nivel de Servicio (SLA), contratos que definen las métricas y tiempos máximos de respuesta esperados por el cliente \cite{Arias_Pro_2025}. En entornos corporativos, la dependencia de procesos de monitoreo manuales y la falta de correlación de eventos generan retrasos inaceptables; esto no solo ocasiona el incumplimiento de los SLA, sino que representa una amenaza directa para la continuidad del negocio y la posición competitiva de la organización \cite{Arias_Pro_2025}.

\textbf{Visualización y Entornos Multi-Nube} \\ 
La gestión de arquitecturas multi-nube (AWS, Azure, GCP) exige estrategias de observabilidad estrictas para evitar puntos ciegos operativos. En esta línea, estudios recientes demuestran que la implementación de herramientas integrales de observabilidad permite identificar cuellos de botella y prevenir la saturación de recursos mediante alertas centralizadas \cite{saavedra_implementacion_2023}. Mientras que dichas implementaciones suelen limitarse a casos de uso aislados o de un solo sector \cite{saavedra_implementacion_2023}, el \textbf{\NombreProyecto} consolida la telemetría de múltiples inquilinos corporativos simultáneamente mediante paneles dinámicos en Grafana y Prometheus, facilitando la toma de decisiones.

\textbf{Middleware e Integración de Sistemas (TIBCO)} \\ 
En entornos corporativos donde conviven múltiples aplicaciones heterogéneas, es necesario un software intermediario ("Middleware"). Soluciones como TIBCO Scribe y TIBCO Business Works permiten automatizar la transferencia de datos y ejecutar catálogos nocturnos, por lo que el monitoreo de sus agentes y la lectura de logs son vitales para la integridad de la información institucional.

\subsection{Tabla comparativa}
\label{sec:tabla_comparativa}

Para sintetizar el análisis del estado del arte, en la Tabla \ref{tab:comparativa_trabajos} se presenta una matriz comparativa que contrasta las características técnicas y operativas del \textbf{\NombreProyecto} frente a los trabajos relacionados descritos anteriormente (Arias, 2025; Saavedra, 2023). Este análisis permite evidenciar cómo la presente propuesta supera las limitantes teóricas e integraciones aisladas al unificar el monitoreo multi-nube, la validación de middleware y la gestión estricta de SLA en un entorno de producción corporativo real.

\begin{table}[H]
\centering
\caption{Comparativa técnica del estado del arte frente a la propuesta}
\label{tab:comparativa_trabajos}
\begin{tabular}{|p{3.2cm}|p{3.6cm}|p{3.6cm}|p{4.2cm}|}
\hline
\textbf{Criterio Evaluado} & \textbf{Proyecto SIMAI (Arias, 2025)} & \textbf{Observabilidad Cloud (Saavedra, 2023)} & \textbf{\NombreProyecto (Propuesta)} \\
\hline
Entorno de Despliegue & Diseño conceptual; simulación de arquitectura en laboratorio. & Entorno de producción real para una sola aplicación financiera. & Entorno de producción multi-inquilino para clientes corporativos. \\
\hline
Alcance de Infraestructura & Homogéneo, limitado nativamente a Microsoft Azure. & Infraestructura Cloud genérica centralizada. & Ecosistema Híbrido y Multi-Nube (AWS, Azure, GCP y On-Premise). \\
\hline
Integración de Ciclo ITSM & Fundamento teórico estructurado bajo ITIL v4. & Fuera del alcance del documento técnico. & Gestión del ciclo de vida y control estricto de SLA (ServiceNow, ZenDesk). \\
\hline
Validación de Middleware (EAI) & No contemplado en la propuesta. & No contemplado en la propuesta. & Supervisión activa de nodos y catálogos en TIBCO Scribe y Business Works. \\
\hline
Observabilidad y Telemetría & Dependencia de Azure Monitor (sin implementación). & Uso de APM comercial dedicado (New Relic). & Consolidación unificada de métricas y depuración (Grafana, Prometheus, DataDog). \\
\hline
\end{tabular}
\end{table}

\subsection{Definición del Problema}
\label{sec:problema}

El problema que se busca resolver es la ineficiencia y el riesgo de interrupción operativa en la gestión de infraestructuras tecnológicas complejas que abarcan entornos locales (On-Premise) y multi-nube. En la actualidad, muchas mesas de servicio operan bajo un modelo reactivo, esperando el reporte de los usuarios en lugar de detectar anomalías preventivamente. Esta falta de monitoreo proactivo, sumada a la saturación de herramientas con alertas redundantes y falsos positivos, genera retrasos significativos en la categorización, escalamiento y resolución de incidentes, lo que frecuentemente resulta en el incumplimiento de los Acuerdos de Nivel de Servicio (SLA) y afecta la continuidad de los aplicativos críticos de los clientes.

Frente a esta problemática, el presente proyecto propone la optimización integral de las operaciones de monitoreo y soporte de TI dentro de un Centro de Operaciones (NOC/SOC). La solución contempla la validación constante de integraciones clave mediante TIBCO (Scribe y Business Works), la gestión rigurosa del ciclo de vida de los incidentes a través de plataformas ITSM (ServiceNow, ZenDesk) y la estructuración de paneles de visualización (Grafana, Prometheus) que permitan consolidar las alertas reales para facilitar la toma de decisiones.

Con base en lo anterior, la pregunta de investigación que guía este trabajo es: ¿De qué manera la optimización del monitoreo proactivo y la gestión estructurada de incidentes ITSM pueden asegurar la alta disponibilidad operativa y reducir el incumplimiento de los SLA en infraestructuras tecnológicas de entornos híbridos y multi-nube?

\subsection{Objetivos}
\label{sec:objetivos}

\subsubsection{Objetivo General}
\label{subsec:objetivo_general}

Asegurar la alta disponibilidad y el rendimiento óptimo de la infraestructura tecnológica y los aplicativos de los clientes, mediante la ejecución de rondines de monitoreo proactivo, la validación de integraciones de software y la gestión eficiente de incidentes. Se espera optimizar la visualización de datos (mediante Grafana y Prometheus) y agilizar la categorización, escalamiento y resolución de tickets en plataformas ITSM, garantizando el cumplimiento de los SLA y mejorando la comunicación entre el nivel operativo y la torre de especialistas (SysAdmin).

\subsubsection{Objetivos Específicos}
\label{subsec:objetivos_especificos}

\begin{itemize}
    \item \textbf{Gestión de Incidentes (ITSM):} Controlar el ciclo de vida de las alertas operativas mediante la creación, asignación, categorización y seguimiento de tickets en plataformas como ServiceNow y ZenDesk, mitigando riesgos de vencimiento en los tiempos de resolución establecidos por los SLA.
    \item \textbf{Supervisión de Integraciones y Aplicativos:} Validar la correcta ejecución de los Manuales de Operación (MOP) correspondientes a las plataformas de integración TIBCO Scribe y TIBCO Business Works, supervisando la estabilidad de los catálogos nocturnos, el reinicio de agentes locales y el análisis de logs de errores en Google Cloud Platform (GCP).
    \item \textbf{Monitoreo Proactivo de Infraestructura:} Identificar de forma temprana anomalías de red, saturación de recursos críticos (CPU, memoria y disco) y estados de inactividad de servidores utilizando herramientas de supervisión como OpManager, AppManager, DataDog y consolas nativas de entornos multi-nube (AWS, Azure y GCP).
    \item \textbf{Optimización de Plataformas de Visualización:} Estructurar y depurar paneles de control centralizados en Grafana y Prometheus mediante el uso de transformaciones y extensiones de tablas dinámicas, consolidando el historial y estado de las alarmas activas con sus respectivas marcas de tiempo para agilizar la toma de decisiones.
\end{itemize}


\subsection{Justificación}
\label{sec:justificacion}

El desarrollo de este proyecto se justifica desde diversas perspectivas debido al impacto directo que tiene en la eficiencia operativa de la organización y en la calidad del servicio entregado a los clientes corporativos.

Desde una perspectiva tecnológica, el proyecto representa una modernización indispensable en la administración de infraestructuras complejas. La transición hacia entornos híbridos y multi-nube exige el dominio de herramientas avanzadas de supervisión y la centralización de datos dispersos. Optimizar el uso de plataformas como Grafana, Prometheus, DataDog y OpManager permite consolidar métricas en tiempo real, transformando datos aislados en paneles de control analíticos que elevan las capacidades de diagnóstico técnico.

Desde el punto de vista operacional y práctico, la relevancia radica en la estandarización y optimización de los procesos de soporte. La validación rigurosa de los Manuales de Operación (MOP) para las interfaces de TIBCO (Scribe y Business Works) asegura la continuidad de catálogos críticos y transferencias de datos sin requerir intervención manual constante. Asimismo, el establecimiento de filtros precisos y ventanas de validación para descartar falsos positivos optimiza el tiempo de la torre de especialistas (SysAdmin), permitiéndoles enfocarse únicamente en incidentes reales que pongan en riesgo la disponibilidad de los aplicativos.

Finalmente, desde un enfoque económico y de negocio, el proyecto mitiga directamente los riesgos financieros asociados a la interrupción de servicios de TI. Las caídas de infraestructura se traducen en severas penalizaciones por el incumplimiento de los Acuerdos de Nivel de Servicio (SLA). Implementar una gestión estructurada y proactiva de incidentes en plataformas ITSM como ServiceNow y ZenDesk garantiza que las alertas se atiendan y documenten dentro de las ventanas de tiempo establecidas, protegiendo la continuidad de operaciones de cuentas de alto perfil y asegurando la calidad del servicio de forma sostenida.

\subsection{Alcances y Limitaciones}
\label{sec:alcances_limitaciones}

\subsubsection{Alcances}
\label{subsec:alcances}

\begin{itemize}
    \item \textbf{Supervisión Multi-Nube e Infraestructura:} El monitoreo abarca tanto servidores locales (\textit{On-Premise}) como instancias en la nube pública a través de las consolas de Amazon Web Services (AWS), Microsoft Azure y Google Cloud Platform (GCP), utilizando herramientas centralizadas como OpManager, AppManager y DataDog.
    \item \textbf{Gestión de Incidentes bajo el Ciclo ITSM:} Cobertura total del flujo de alertas mediante el uso de ServiceNow y ZenDesk para el registro, categorización, asignación y seguimiento de tickets, implementando un esquema de documentación cruzada (folios de notificación para el cliente y folios de análisis interno).
    \item \textbf{Validación de Integraciones Middleware:} Supervisión del estado y ejecución de las interfaces críticas mediante TIBCO Scribe y TIBCO Business Works, incluyendo el control de los catálogos nocturnos, el reinicio guiado de agentes y la revisión de logs de Error Reporting en GCP.
    \item \textbf{Optimización y QA en Plataformas de Monitoreo:} Creación de entornos de prueba controlados para la depuración y ajuste de consultas (\textit{queries}) en Grafana y Prometheus, implementando complementos dinámicos como \textit{Business Table} para enriquecer la inserción de comentarios técnicos.
\end{itemize}

\subsubsection{Limitaciones}
\label{subsec:limitaciones}

\begin{itemize}
    \item \textbf{Dependencia de Licenciamiento y Permisos:} El acceso a las métricas avanzadas, consolas de infraestructura corporativa y plataformas de mesa de servicio está restringido por los privilegios de usuario y contratos de licenciamiento previamente establecidos por la organización y sus clientes de terceros.
    \item \textbf{Modalidad de Trabajo Híbrida:} La ejecución del proyecto se realiza bajo un esquema mixto (\textit{Home Office} y días presenciales). Esto limita la comunicación directa interpersonal, haciendo que el escalamiento de incidentes hacia los SysAdmin y la resolución de alertas dependan críticamente de la disponibilidad de herramientas de mensajería (Teams, WhatsApp) y de la estabilidad en los accesos remotos por VPN.
    \item \textbf{Calidad de los Datos de Entrada:} La eficiencia en el filtrado de alarmas y estados de error en Grafana está condicionada por la estructura y formato de los registros históricos de logs generados originalmente por los sistemas del cliente, no teniendo control sobre anomalías previas en el código fuente de los aplicativos.
\end{itemize}

\subsection{Temas tentativos del marco teórico}
\label{sec:temas_tentativos_marco_teorico}

\begin{itemize}
    \item Gestión de Servicios de TI (ITSM) y el ciclo de vida de los incidentes.
    \item Acuerdos de Nivel de Servicio (SLA) y su impacto en la continuidad operativa.
    \item Centros de Operaciones (NOC/SOC) y la transición del soporte reactivo al monitoreo proactivo.
    \item Arquitecturas de infraestructura tecnológica en entornos Híbridos y Multi-Nube (AWS, Azure, GCP).
    \item Middleware e Integración de Aplicaciones Empresariales (EAI) utilizando ecosistemas TIBCO.
    \item Telemetría, observabilidad y herramientas de monitoreo de red (OpManager, DataDog, AppManager).
    \item Visualización de datos y consolidación de métricas operativas (Grafana y Prometheus).
    \item Estrategias de reducción de ruido, filtrado de alertas y depuración de falsos positivos.
\end{itemize}

\iffalse
\subsection{Ejemplos de Tablas e Inclusión de Imágenes}
\label{sec:ejemplos_elementos}

Para ilustrar el uso de elementos visuales en su tesis, a continuación se presentan ejemplos básicos de cómo incluir tablas e imágenes en LaTeX.

\subsubsection{Inclusión de Tablas}

Las tablas son herramientas fundamentales para presentar datos de manera organizada y comprensible.
\begin{table}[H]
    \centering
    \caption{Clasificación de Sensores Comunes en Ingeniería.}
    \label{tab:sensores_comunes}
    \begin{tabular}{|l|c|c|}
        \hline
        \textbf{Tipo de Sensor} & \textbf{Magnitud Medida} & \textbf{Principio de Funcionamiento} \\
        \hline
        Termopar & Temperatura & Efecto Seebeck \\
        Extensómetro & Deformación & Resistencia eléctrica \\
        Acelerómetro & Aceleración & Efecto piezoeléctrico \\
        Encoder rotatorio & Posición angular & Óptico/Magnético \\
        \hline
    \end{tabular}
    \caption*{Nota: Esta tabla es un ejemplo para demostrar la estructura básica de una tabla en LaTeX.}
\end{table}
Como se puede observar en la Tabla \ref{tab:sensores_comunes}, la información se presenta de forma clara y estructurada. Es importante añadir un \verb|`\caption`|  para describir la tabla y un  \verb|`\label`|  para referenciarla en el texto.

\subsubsection{Inclusión de Imágenes}

Las imágenes, gráficos o diagramas son cruciales para complementar la explicación textual y facilitar la comprensión de conceptos complejos o resultados visuales.
\begin{figure}[H]
    \centering
    \includegraphics[width=0.7\textwidth]{img/diagrama_proceso}
    \caption{Diagrama de un proceso de control automático.}
    \label{fig:diagrama_control}
\end{figure}
% La Figura \ref{fig:diagrama_control} ilustra un ejemplo de un diagrama de un proceso de control. Para incluir una imagen, se utiliza el entorno `figure` y el comando `\includegraphics`. Es fundamental especificar la ruta de la imagen (en este caso, `images/diagrama_proceso.png`) y un `width` para controlar su tamaño. Al igual que con las tablas, un `\caption` y un `\label` son indispensables. Es crucial que cada imagen tenga una calidad adecuada y que sea relevante para el texto.

\begin{figure}[H]
    \centering
    \begin{subfigure}[b]{0.45\textwidth}
        \includegraphics[width=\textwidth]{img/grafico_datos1}
        \caption{Datos de Temperatura.}
        \label{fig:subfig_temp}
    \end{subfigure}
    \hfill
    \begin{subfigure}[b]{0.45\textwidth}
        \includegraphics[width=\textwidth]{img/grafico_datos2}
        \caption{Datos de Presión.}
        \label{fig:subfig_pres}
    \end{subfigure}
    \caption{Comparación de datos operativos de un sistema (Temperatura vs. Presión).}
    \label{fig:comparacion_datos}
\end{figure}
Incluso es posible combinar varias imágenes en una sola figura, como se muestra en la Figura \ref{fig:comparacion_datos}, utilizando el paquete `subcaption`.

Asegúrense de que todas las figuras y tablas estén debidamente referenciadas en el texto y que su contenido sea autoexplicativo, utilizando los pies de figura y títulos de tabla de manera efectiva.

\fi